\begin{pastexam}[label={px:ncu-112-dsa-05}]{中央資工所 112 學年度 複選題第 5 題}
（複選；每題 5 分，答錯 1 題倒扣 1 分）

Which of the following statement(s) is (are) correct?
\begin{enumerate}[label=(\Alph*)]
\item The time complexity of inserting an element to an n-node min-Heap is O(n).
\item The time complexity of searching an element in an n-node max-Heap is O(n).
\item Building the min-Heap by inserting these nodes one by one \{42, 9, 23, 37, 4, 34, 2\}. The level
order traversal sequence of the min-Heap is \{2, 9, 4, 42, 37, 23, 34\}.
\item The time complexity of deleting min in an n-node min-Heap is O(log n)
\end{enumerate}
\end{pastexam}

\begin{pastsoln}
答案是 (B)、(D)。

這類敘述題的 $O(\cdot)$ 指的是該操作的\textbf{緊的}複雜度（同一題 (B)、(D) 也是這樣用）。 (A) 若只當上界
看，$O(\log n) \subseteq O(n)$，字面上也成立；但出題的意思顯然是「插入要 $O(n)$」，下面照這個意思判斷。

\begin{itemize}
\item (A) 錯。 插入是把新元素放在最後一格再往上調整，最多走樹高 $\lfloor \lg n \rfloor$ 層，
$O(\log n)$（資結書第~8.1~節、虛擬碼~8.2）。
\item (B) 對。 最大堆積只保證父節點不小於子節點，同一層、左右子樹之間沒有大小順序，找一個任意的值
無法像二元搜尋樹那樣排除一半。 最壞情況（例如要找的值不在堆積裡，或比較小而藏在某片葉子）要看遍
所有節點，$O(n)$。
\item (C) 錯。 逐一插入、每次往上調整（陣列從 $1$ 開始，父節點是 $\lfloor i/2 \rfloor$）：
\begin{center}
\small
\begin{tabular}{lll}
\toprule
插入 & 往上調整 & 陣列（層序） \\
\midrule
42 & & 42 \\
9 & $9 < 42$，交換 & 9, 42 \\
23 & $23 > 9$，不動 & 9, 42, 23 \\
37 & 放第 $4$ 格，父是 $42$，交換；再和 $9$ 比，不動 & 9, 37, 23, 42 \\
4 & 放第 $5$ 格，父是 $37$，交換；再和 $9$ 比，交換 & 4, 9, 23, 42, 37 \\
34 & 放第 $6$ 格，父是 $23$，不動 & 4, 9, 23, 42, 37, 34 \\
2 & 放第 $7$ 格，父是 $23$，交換；再和 $4$ 比，交換 & 2, 9, 4, 42, 37, 34, 23 \\
\bottomrule
\end{tabular}
\end{center}
層序是 $\{2, 9, 4, 42, 37, 34, 23\}$，最後兩個和題目的 $\{\ldots, 23, 34\}$ 順序相反：$2$ 從第 $7$ 格
（$23$ 的右子節點）往上換，$23$ 被換到第 $7$ 格，$34$ 一直留在第 $6$ 格。
\item (D) 對。 刪除最小值是拿走根，把最後一格搬到根再往下調整，最多走樹高，$O(\log n)$。
\end{itemize}
\end{pastsoln}
